\chapter{INTRODUCTION}

\section{Background of the Study}

Medical image segmentation plays a crucial role in clinical applications, particularly in the context of brain tumor detection and diagnosis using Magnetic Resonance Imaging (MRI) \parencite{cai2022novel, liu2021review}. Accurate segmentation is vital for identifying tumor boundaries, assessing tumor characteristics, and guiding treatment decisions. This process separates images into different regions with similar characteristics, aiding in the intuitive interpretation of medical images and reducing errors \parencite{liu2021review, handels2013viewpoints}.

The broader research domain of computer-aided diagnosis (CAD) has been revolutionized by deep learning techniques, specifically Convolutional Neural Networks (CNNs), which have proven to be more accurate and robust than traditional methods for medical image segmentation. Deep learning architectures like U-Net and its variants have been extensively used for organ and lesion segmentation in various medical imaging modalities, including MRI scans for brain tumor analysis \parencite{cai2022novel, moorthy2022survey, tang2019role}.

Accurate and efficient brain tumor segmentation is critical for improving patient outcomes through early detection, precise treatment planning, and monitoring tumor progression over time \parencite{liu2021review, tang2019role}. Automated segmentation methods based on deep learning not only enhance the speed and accuracy of tumor delineation but also assist in reducing the burden on radiologists and healthcare providers \parencite{moorthy2022survey, tang2019role}.

Despite the advancements in deep learning-based brain tumor segmentation approaches, there are limitations and challenges that need to be addressed. The recent study by \textcite{hu2022net} on the "AS-Net: Attention Synergy Network for skin lesion segmentation" has shown the effectiveness of using attention mechanisms, such as spatial and channel attention, for improving image segmentation performance. However, the authors of the said study acknowledged limitations in the generalizability and computational efficiency of their model due to the use of a relatively complex encoder architecture, namely VGG16.

Inspired by the success of the AS-Net model in skin lesion segmentation, the proposed thesis idea aims to adapt and optimize the AS-Net approach for MRI-based brain tumor segmentation. Specifically, this study will investigate the performance of the AS-Net architecture with alternative, lightweight encoder networks, such as MobileNetV3 and EfficientNetV2, in comparison to the original VGG16 encoder. 

MobileNetV3 and EfficientNetV2 are chosen as alternative encoders due to their demonstrated ability to achieve high accuracy with significantly reduced computational cost and memory footprint. These lightweight architectures are well-suited for deployment on resource-constrained devices or in scenarios where real-time performance is crucial, making them particularly relevant for medical image analysis applications.

The goal is to achieve improved segmentation accuracy while enhancing the computational efficiency of the model, ultimately contributing to the development of more practical and clinically-relevant brain tumor segmentation solutions.


\section{Statement of the Problem}

The recent study on the "AS-Net: Attention Synergy Network for skin lesion segmentation" \textcite{hu2022net} has demonstrated the effectiveness of attention mechanisms in enhancing image segmentation performance. However, the previous studies regarding the AS-Net acknowledged limitations in the generalizability and computational efficiency of their model due to the use of the computationally heavy VGG16 encoder architecture.

Specifically, this research problem seeks to address the following:

\begin{enumerate}
    \item By exploring alternative lightweight models, how can MobileNetV3 and EfficientNetV2 be effectively integrated into the existing AS-Net architecture?
    \item How efficient are the encoders of MobileNetV3 and EfficientNetV2 in terms of processing speed and performance (measured by metrics like F1 score, Jaccard index, precision-recall curve, confusion matrix, ROC AUC score, and ROC curve) compared to VGG16?
\end{enumerate}


\section{Objectives of the Study}

The proposed thesis research aims to achieve the following objectives:

Specifically, this study aims to:

\begin{enumerate}
    \item Primary objective:
    \begin{itemize}
        \item To optimize the AS-Net architecture for MRI-based brain tumor segmentation by comparing the performance of different encoder networks, including lightweight models.
    \end{itemize}
    \item Secondary objectives:
    \begin{itemize}
        \item Evaluate the effectiveness of spatial and channel attention mechanisms for brain tumor segmentation in MRI images.
        \item Assess the segmentation performance, computational efficiency, and generalizability of AS-Net with lightweight encoder architectures (e.g., MobileNetV3, EfficientNetV2) compared to VGG16.
        \item Provide recommendations for the most suitable encoder network to be integrated into the AS-Net architecture for improved MRI-based brain tumor segmentation.
    \end{itemize}
\end{enumerate}


\section{Significance of the Study}

This proposed research holds significant implications for the advancement of knowledge in the field of medical image analysis and computer-aided diagnosis (CAD) for brain tumors. The adaptation and optimization of the AS-Net architecture \parencite{hu2022net} for MRI-based brain tumor segmentation will contribute to the growing body of literature on deep learning-based solutions for this critical medical imaging task.

The exploration of lightweight encoder architectures, such as MobileNetV3 and EfficientNetV2, within the AS-Net framework has the potential to enhance the practical applicability of deep learning-based brain tumor segmentation. By achieving improved computational efficiency without compromising segmentation accuracy, the proposed model can facilitate the deployment of these advanced techniques in real-world clinical settings, where computational resources and processing time are often limited.

The findings of this study will be particularly beneficial for various stakeholders in the healthcare ecosystem:

\begin{enumerate}
    \item Clinicians and radiologists involved in the diagnosis and management of brain tumors will have access to a more accurate, efficient, and accessible brain tumor segmentation tool, which can significantly enhance their decision-making capabilities and improve patient outcomes.
    \item Patients will benefit from earlier and more reliable detection of brain tumors, leading to timely interventions and potentially better prognosis.
\end{enumerate}

Furthermore, the adaptation of the attention synergy mechanism, a key component of the original AS-Net model, to the domain of brain tumor segmentation will contribute to the understanding of the role of attention-based feature extraction in medical image analysis. This knowledge can be leveraged to further refine and optimize deep learning architectures for a wide range of medical imaging applications, ultimately advancing the field of computer-aided diagnosis and accelerating the translation of these technologies into clinical practice


\section{Scope and Limitation of the Study}

This proposed research focuses on the adaptation and optimization of the AS-Net architecture for the task of MRI-based brain tumor segmentation. The study will primarily concentrate on the comparison of different encoder architectures, namely VGG16, MobileNetV3, and EfficientNetV2, within the AS-Net framework. By evaluating the performance, computational efficiency, and generalizability of these encoder networks, the study aims to provide recommendations for the most suitable encoder to be integrated into the AS-Net model for improved brain tumor segmentation.

While the scope of this research is well-defined, there are certain limitations that should be acknowledged. The study will be conducted using publicly available dataset for MRI-based brain tumor segmentation. The characteristics of the chosen dataset, including the diversity and distribution of brain tumor types, may influence the generalizability of the research findings. As such, the performance of the optimized AS-Net model may need to be further validated using additional brain tumor segmentation datasets.

Moreover, the availability and quality of the MRI data used in the study could pose challenges during the research process. Factors such as image resolution, slice thickness, and the presence of artifacts or noise may affect the model's performance and require careful data preprocessing and augmentation techniques. Additionally, the training and optimization of the deep learning models may be constrained by the computational resources available, necessitating the exploration of efficient training strategies and model compression techniques.

Despite these limitations, the proposed research holds significant potential to contribute to the advancement of deep learning-based brain tumor segmentation. By exploring the performance of lightweight encoder architectures within the AS-Net framework, the study aims to provide practical and clinically-relevant solutions that can be readily deployed in real-world healthcare settings, ultimately enhancing the accuracy, efficiency, and accessibility of brain tumor detection and monitoring.